% !TeX program = xelatex
% ============================================================================
%  第 4 课　群的四条规则
% ----------------------------------------------------------------------------
%  编译        XeLaTeX，编译两遍
%  样式        ../../00-共享/preamble.tex
%  说明        本课不发单页
%  图片        一律 \imgph 占位，不要在这里手画 tikz
%  记号        新定义运算统一写作 \oplus（避免在数学模式里排中文符号）
%  说明        「运算」的定义已在第 2 课讲过，本课不再重复，只做一句回顾
% ============================================================================

\documentclass[10pt]{beamer}
% ============================================================================
%  讲义样式文件 —— 字体 / 配色 / 校徽栏 / 顶部导航栏 / 页脚 / 常用方框
%  所有个性化设置都集中在这里，改一次即可用于所有讲义
%
%  用法：主文件里写
%      \documentclass[10pt]{beamer}
%      \input{preamble.tex}
%  依赖：同目录下的 logo_white.png（顶部校徽）
%  编译：必须用 XeLaTeX，而且必须编译两遍
%        （顶部导航栏的节名与小圆点记录在 .nav 文件里，只编一遍不会更新）
% ============================================================================
\usepackage{ctex}  % 支持中文
\usepackage{amsmath, amssymb}  % 数学公式支持
\usepackage{fontspec}  % 字体设置
\usepackage{unicode-math}  % 支持Unicode数学字体
\usepackage{xcolor}  % 颜色
\usepackage{tikz}  % 用于绘制页脚色条
\usepackage{booktabs}  % 表格

% 图片搜索路径 ------------------------------------------------------------------
% 校本课的资料分目录存放，本文件可能被 02-课件/第XX课/ 、03-讲义/ 下的主文件调用，
% 这里把几种可能的相对位置都列上，logo_white.png 放在 00-共享/ 即可被找到。
% 最后一项 00-共享/图片/ 是给 \imgfig 用的插图目录（图与课件不在同一层，靠它找到）。
\graphicspath{{./}{./图片/}{../../00-共享/}{../00-共享/}{00-共享/}{../../00-共享/图片/}{../00-共享/图片/}{00-共享/图片/}}

% 设置字体
\setCJKmainfont{SimSun}  % 设置中文字体为宋体
\setmainfont{Palatino Linotype}  % 设置英文字体为Palatino Linotype
\setmathfont{Latin Modern Math}  % 设置数学字体为Latin Modern Math
% 注意：beamer 默认把正文字体族设成无衬线（sans），Windows 下对应的中文字体是
% 微软雅黑，所以上面这行 SimSun 目前对正文不起作用。想让正文变成宋体，取消
% 下面这行的注释即可（标题栏、页脚、方框标题不受影响）。
% \renewcommand{\familydefault}{\rmdefault}

% 定义颜色 - 均为绿色系
\definecolor{maingreen}{RGB}{34,139,34}  % 森林绿（导航栏底色）
\definecolor{lightgreen}{RGB}{240,255,240}  % 蜜瓜白（帧标题栏底色）
\definecolor{accentgreen}{RGB}{46,139,87}  % 海绿（强调色）
\definecolor{textgreen}{RGB}{0,128,0}  % 纯绿（标题文字）
\definecolor{darkgreen}{RGB}{20,90,20}  % 深绿（顶部校徽栏底色）

% 设置Beamer主题
\usetheme{default}
\usecolortheme{default}

% 自定义主题
\setbeamercolor{background canvas}{bg=white}
\setbeamercolor{frametitle}{fg=textgreen, bg=lightgreen}
\setbeamercolor{title}{fg=textgreen}
\setbeamercolor{subtitle}{fg=textgreen}
\setbeamercolor{section in toc}{fg=accentgreen}
\setbeamercolor{normal text}{fg=black}
\setbeamercolor{structure}{fg=accentgreen}
\setbeamercolor{itemize item}{fg=accentgreen}
\setbeamercolor{itemize subitem}{fg=accentgreen}
\setbeamercolor{itemize subsubitem}{fg=accentgreen}
\setbeamercolor{block title}{fg=white, bg=accentgreen}
\setbeamercolor{block body}{bg=lightgreen}
% example 环境（底层是 exampleblock）用的是独立的一组颜色名，default 配色主题下
% 它们的底色是空的，所以只设 block 的话“例”会变成一行没有框的标题，这里接过来
\setbeamercolor{block title example}{parent=block title}
\setbeamercolor{block body example}{parent=block body}

% ==================== 常用定理类方框（五类配色）====================
% 五类各自的标题条底色，正文底色由同色自动调浅（!8!white）
\definecolor{boxThm}{RGB}{20,90,20}    % 一、定理类：深绿
\definecolor{boxDef}{RGB}{46,139,87}   % 二、定义类：海绿
\definecolor{boxExm}{RGB}{34,139,34}   % 三、例题类：森林绿
\definecolor{boxPrf}{RGB}{85,107,47}   % 四、过程类：橄榄绿
\definecolor{boxNote}{RGB}{0,121,107}  % 五、注记类：青绿

% 方框本体：先套本类配色，再调用 beamer 的 block（改动只在环境内部生效）
\newcommand{\gbox}[2]{%
  \setbeamercolor{block title}{fg=white,bg=#1}%
  \setbeamercolor{block body}{bg=#1!8!white}%
  \begin{block}{#2}}
\newcommand{\gboxend}{\end{block}}
% 标题后的可选补充名：\begin{theorem}[勾股定理] → 定理（勾股定理）
\newcommand{\gsub}[1]{\ifx\relax#1\relax\else（#1）\fi}


% 一、定理类（定理/推论/性质/命题/引理）
\renewenvironment{theorem}[1][]{\gbox{boxThm}{定理\gsub{#1}}}{\gboxend}
\renewenvironment{corollary}[1][]{\gbox{boxThm}{推论\gsub{#1}}}{\gboxend}
\newenvironment{character}[1][]{\gbox{boxThm}{性质\gsub{#1}}}{\gboxend}
\newenvironment{proposition}[1][]{\gbox{boxThm}{命题\gsub{#1}}}{\gboxend}
\renewenvironment{lemma}[1][]{\gbox{boxThm}{引理\gsub{#1}}}{\gboxend}

% 二、定义类（定义/公理）
\renewenvironment{definition}[1][]{\gbox{boxDef}{定义\gsub{#1}}}{\gboxend}
\newenvironment{axiom}[1][]{\gbox{boxDef}{公理\gsub{#1}}}{\gboxend}

% 三、例题类（例/练习/习题）
\renewenvironment{example}[1][]{\gbox{boxExm}{例\gsub{#1}}}{\gboxend}
\newenvironment{exercise}[1][]{\gbox{boxExm}{练习\gsub{#1}}}{\gboxend}
\newenvironment{homework}[1][]{\gbox{boxExm}{习题\gsub{#1}}}{\gboxend}

% 四、过程类（证明/解/法1~法4）
\renewenvironment{proof}[1][]{\gbox{boxPrf}{证明\gsub{#1}}}{\qedsymbol\gboxend}
\newenvironment{jie}[1][]{\gbox{boxPrf}{解\gsub{#1}}}{\gboxend}
\newenvironment{way1}[1][]{\gbox{boxPrf}{法1\gsub{#1}}}{\gboxend}
\newenvironment{way2}[1][]{\gbox{boxPrf}{法2\gsub{#1}}}{\gboxend}
\newenvironment{way3}[1][]{\gbox{boxPrf}{法3\gsub{#1}}}{\gboxend}
\newenvironment{way4}[1][]{\gbox{boxPrf}{法4\gsub{#1}}}{\gboxend}

% 五、注记类（经验总结/注/值得注意的是/分析/提示）
\newenvironment{experience}[1][]{\gbox{boxNote}{经验总结\gsub{#1}}}{\gboxend}
\newenvironment{zhu}[1][]{\gbox{boxNote}{注\gsub{#1}}}{\gboxend}
\newenvironment{att}[1][]{\gbox{boxNote}{值得注意的是\gsub{#1}}}{\gboxend}
\newenvironment{analyse}[1][]{\gbox{boxNote}{分析\gsub{#1}}}{\gboxend}
\newenvironment{hint}[1][]{\gbox{boxNote}{提示\gsub{#1}}}{\gboxend}

% 顶部导航栏（miniframes 主题：显示各节名称 + 帧圆点）
% 节名与圆点记录在 .nav 文件中，改动章节后需要编译两遍才会同步
\useoutertheme[subsection=false]{miniframes}
\setbeamercolor{section in head/foot}{fg=yellow, bg=maingreen}
\setbeamercolor{mini frame}{fg=yellow, bg=maingreen}
\setbeamercolor{mini frame in current subsection}{fg=yellow, bg=accentgreen}

% 顶部校徽栏：位于导航栏上方，宽度与帧标题栏一致（整页宽）
% 校训用的楷体：ctex 的 \kaishu 在本机指向未安装的 Adobe 楷体，这里直接指定系统楷体
\setCJKfamilyfont{kaiti}{KaiTi}
\setbeamercolor{logo bar}{fg=white, bg=darkgreen}
\addtobeamertemplate{headline}{%
  \begin{beamercolorbox}[wd=\paperwidth,ht=0.84cm,dp=0.20cm,leftskip=0.25cm,rightskip=0.25cm]{logo bar}%
    % -0.035cm：色条内容默认坐在基线上，会比栏的中线略高，下移这一点使其上下居中
    \raisebox{-0.035cm}{\includegraphics[height=0.72cm]{logo_white.png}}%
    % 右侧校训：字体为上面定义的楷体，也可改为 \insertsectionhead（当前节名）、\inserttitle（章节名）等
    % 0.375cm = 校徽中线（0.36cm）加汉字基线偏移，并同样下移 0.035cm 与校徽对齐
    \hfill\raisebox{\dimexpr 0.375cm - 0.5\height - 0.5\depth\relax}%
      {\CJKfamily{kaiti}\fontsize{12pt}{24pt}\selectfont 自强不息，止于至善}%
  \end{beamercolorbox}%
}{}

% % 设置边框 - 使用TikZ绘制圆角
% \setbeamertemplate{frametitle}{
%   \vspace*{0.2cm}
%   \begin{tikzpicture}
%     % 背景矩形宽度限制为标题内容宽度
%     \fill[lightgreen, rounded corners=5pt] (0,0) rectangle (\textwidth+0.6cm, 0.8cm);
%     \node[anchor=west, text width=\textwidth, xshift=0.3cm, yshift=0.4cm] {\usebeamerfont{frametitle}\insertframetitle};
%   \end{tikzpicture}
% }

% 设置项目符号
\setbeamertemplate{itemize item}{\small\raise0.5pt\hbox{\textcolor{accentgreen}{$\bullet$}}}
\setbeamertemplate{itemize subitem}{\small\raise0.3pt\hbox{\textcolor{accentgreen}{$\circ$}}}

% 设置页脚 - 使用TikZ绘制一条整页宽的色条，左侧页码、右侧姓名
\setbeamertemplate{footline}{
  \begin{tikzpicture}
    \fill[maingreen] (-1,0) rectangle (\paperwidth, 0.4cm);
    \node[yellow,anchor=east, xshift=0cm, yshift=0.2cm] {\insertframenumber/\inserttotalframenumber};
    % 原点在纸张左边缘右侧 1cm 处（色条从 x=-1cm 画起），故 \paperwidth-2cm
    % 对应纸张右边缘内缩 1cm 的位置，与左侧页码左右对称
    \node[yellow,anchor=east, xshift=\paperwidth-2cm, yshift=0.2cm] {\insertauthor};
  \end{tikzpicture}
}

% 关闭 beamer 默认的导航符号（页码与姓名已在页脚中自行绘制）
\setbeamertemplate{navigation symbols}{}

% 设置标题页面
\setbeamertemplate{title page}{
  \vspace{0.5cm}
  \begin{center}
    \begin{minipage}[c][0.68\paperheight][c]{\textwidth}
      {\usebeamerfont{title}\usebeamercolor[fg]{title}\large\inserttitle\par}
      \vspace{1.5cm}
      {\usebeamerfont{subtitle}\usebeamercolor[fg]{subtitle}\Large\insertsubtitle\par}
      \vspace{1.5cm}
      {\usebeamerfont{author}\usebeamercolor[fg]{author}\insertauthor\par}
      \vspace{0.5cm}
      {\usebeamerfont{institute}\usebeamercolor[fg]{institute}\insertinstitute\par}
      \vspace{0.5cm}
      {\usebeamerfont{date}\usebeamercolor[fg]{date}\insertdate\par}
    \end{minipage}
  \end{center}
}

% 设置块环境
\setbeamertemplate{blocks}[rounded][shadow=false]

% 设置目录
\setbeamertemplate{section in toc}{\inserttocsectionnumber.~\inserttocsection}


% ----------------------------------------------------------------------------
%  图片占位符（临时用；有图以后把那行换成 \includegraphics 即可）
%      \imgph{说明文字}          默认高 3.6cm
%      \imgph[6cm]{说明文字}     指定高度
%  会画一个绿色虚线框，框里写说明，方便对着 PDF 找位置。
\newcommand{\imgph}[2][3.6cm]{%
  \begin{center}%
  \begin{tikzpicture}%
    \node[draw=accentgreen, densely dashed, line width=1pt, rounded corners=4pt,
          fill=lightgreen, inner sep=8pt, align=center, text width=0.72\textwidth,
          minimum height=#1, font=\small]
      {\textcolor{accentgreen}{\bfseries 【图片占位】}\\[4pt] #2};%
  \end{tikzpicture}%
  \end{center}%
}

% ----------------------------------------------------------------------------
%  插图（图已就位后用这个；图放在 00-共享\图片\ 下，由 \graphicspath 找到）
%      \imgfig{文件名}                 默认宽 0.72\textwidth
%      \imgfig[0.55\textwidth]{文件名}  指定宽度
%      \imgfig[8cm]{文件名}            也可以直接给尺寸
%  图源与编译方法见 00-共享\图片\README.md。每张图是一页 PDF，只含图、无页边。
\newcommand{\imgfig}[2][0.72\textwidth]{%
  \begin{center}%
    \includegraphics[width=#1]{#2}%
  \end{center}%
}

% ----------------------------------------------------------------------------
%  可用方框环境速查（五类配色，均不编号，均支持 \begin{xxx}[补充名]）
%    定理类（深绿）  theorem 定理 / corollary 推论 / character 性质 /
%                    proposition 命题 / lemma 引理
%    定义类（海绿）  definition 定义 / axiom 公理
%    例题类（森林绿）example 例 / exercise 练习 / homework 习题
%    过程类（橄榄绿）proof 证明 / jie 解 / way1 法1 / way2 法2 /
%                    way3 法3 / way4 法4
%    注记类（青绿）  experience 经验总结 / zhu 注 / att 值得注意的是 /
%                    analyse 分析 / hint 提示
%  另有 beamer 自带的普通方框：\begin{block}{标题} ... \end{block}
% ----------------------------------------------------------------------------


% ==================== 每份课件只改这一段 ====================
\title{从数数开始的抽象代数生活}
\subtitle{第 4 课　群的四条规则}
\author{姜博瀚}
\institute{厦门大学附属实验中学}
\date{\today}
% ============================================================

\begin{document}

\begin{frame}
  \titlepage
\end{frame}

% ---------------- 本课目标 ----------------
\begin{frame}{本课目标}
  \begin{block}{本课目标}
    \begin{itemize}
      \item 能说出群的四条规则。
      \item 会在模 \(12\) 加法、有理数加法、新定义运算里各找到一个例子。
      \item 知道单位元不一定是 \(0\)，交换律不是群的必要条件，结合律也不是废话。
    \end{itemize}
  \end{block}
\end{frame}

% ============================================================================
\section{回顾}

\begin{frame}{回顾：上节课的三个词}
  \begin{center}
  \begin{tabular}{@{}lcc@{}}
    \toprule
     & \textbf{单位元} & \textbf{逆元} \\
    \midrule
    模 \(12\) 加法 & \(0\) & \(12-a\) \\
    有理数加法 & \(0\) & \(-a\) \\
    有理数乘法 & \(1\) & \(1/a\ (a\neq0)\) \\
    \bottomrule
  \end{tabular}
  \end{center}
  \pause
  \begin{center}
    运算、单位元、逆元——这三个词上一节课都讲过了。

    \pause

    {\large 这三套运算，看上去很像。\pause 像在哪儿？}
  \end{center}
\end{frame}

% ============================================================================
\section{四条规则}


\begin{frame}{群：四条规则}
  \begin{definition}[群]
    设 \(*\) 是集合 \(G\) 上的一个运算。如果下面四条都成立，
    就说 \((G,*)\) 是一个\textbf{群}：
    \begin{enumerate}
      \item \textbf{封闭}：任取 \(a,b\in G\)，\(a*b\) 仍然属于 \(G\)；
      \item \textbf{结合}：\((a*b)*c=a*(b*c)\)；
      \item \textbf{单位元}：存在 \(e\in G\)，对每个 \(a\) 都有 \(a*e=e*a=a\)；
      \item \textbf{逆元}：每个 \(a\in G\) 都有 \(b\in G\)，使 \(a*b=b*a=e\)。
    \end{enumerate}
  \end{definition}
  \pause
  \begin{att}
    看清楚：\textbf{交换律不在里面}。\(a*b=b*a\) 不要求成立。
  \end{att}
\end{frame}

% ============================================================================
\section{四个例子}

\begin{frame}{拿四个东西来对一对}
  \begin{center}
  \begin{tabular}{@{}lcccc@{}}
    \toprule
     & 封闭 & 结合 & 单位元 & 逆元 \\
    \midrule
    模 \(12\) 加法 & \(\checkmark\) & \(\checkmark\) & \(0\) & \(12-a\) \\
    有理数加法 & \(\checkmark\) & \(\checkmark\) & \(0\) & \(-a\) \\
    有理数乘法 & \(\checkmark\) & \(\checkmark\) & \(1\) & \(1/a\ (a\neq0)\) \\
    \(a\oplus b=a+b+1\) & ? & ? & ? & ? \\
    \bottomrule
  \end{tabular}
  \end{center}
  \pause
  \begin{att}
    第三行请特别看一眼：\(a=0\) 的时候，\(1/a\) 是什么？—— 没有。
  \end{att}
\end{frame}

\begin{frame}{用手指回答}
  先在心里想好，再举手：
  \begin{itemize}
    \item 模 \(12\) 加法有单位元吗？—— 有的话举 \(1\) 根手指。
    \item 有理数乘法呢？—— 先想一想，再举。
  \end{itemize}
  \pause
  \begin{block}{答案}
    都有。模 \(12\) 加法的单位元是 \(0\)；有理数乘法的单位元是 \(1\)。\\
    （单位元是跟着\textbf{运算}走的，不是永远都叫 \(0\)。）
  \end{block}
\end{frame}

\begin{frame}{有理数乘法：卡在 0 上}
  \begin{block}{\(0\) 没有倒数}
    要找 \(0\) 的逆元，就是要找一个数 \(b\)，使
    \[
      0\times b=1 .
    \]
    可是左边永远是 \(0\)。这样的 \(b\) 根本不存在。
  \end{block}
  \pause
  \begin{center}
    {\large 少了一个逆元，四条规则就凑不齐。}\\[0.5em]
    \pause
    {\large 所以：\textbf{有理数乘法不构成群}。}
  \end{center}
\end{frame}

\begin{frame}{自己定制一个运算}
  在有理数上规定一个新运算：
  \[
    a\oplus b = a+b+1 .
  \]
  \pause
  \begin{example}
    \(3\oplus5=3+5+1=9\)；\quad \((-2)\oplus0=-2+0+1=-1\)。
  \end{example}
  \pause
  有的题目就爱这样出新定义运算。\pause 今天，我们自己当出题人，一条一条验给它看。
\end{frame}

\begin{frame}{验一验：前两条}
  \begin{block}{封闭}
    \(a\oplus b=a+b+1\) 是两个有理数相加、再加 \(1\)，结果还是有理数。\quad \(\checkmark\)
  \end{block}
  \pause
  \begin{block}{结合律}
    \begin{align*}
      (a\oplus b)\oplus c &= (a+b+1)+c+1 = a+b+c+2,\\
      a\oplus(b\oplus c) &= a+(b+c+1)+1 = a+b+c+2 .
    \end{align*}
    两边一样。\quad \(\checkmark\)
  \end{block}
\end{frame}

\begin{frame}{单位元是谁？}
  设单位元是 \(e\)，它要满足 \(a\oplus e=a\)：
  \[
    a+e+1=a \quad\Longrightarrow\quad e=-1 .
  \]
  \pause
  \begin{center}
    {\large 单位元是 \(\mathbf{-1}\)，\textbf{不是 \(0\)}。}
  \end{center}
  \pause
  \begin{block}{验证}
    \(a\oplus(-1)=a+(-1)+1=a\)，对每个 \(a\) 都对。\quad \(\checkmark\)
  \end{block}
\end{frame}

\begin{frame}{每个数的搭档}
  要找 \(a\) 的逆元，就是解 \(a\oplus b=-1\)：
  \[
    a+b+1=-1 \quad\Longrightarrow\quad b=-a-2 .
  \]
  \pause
  \begin{center}
    {\large \(a\) 的逆元是 \(-a-2\)。}
  \end{center}
\end{frame}

\begin{frame}{用手指回答}
  在 \(a\oplus b=a+b+1\) 里，报出下面每个数的逆元：
  \begin{itemize}
    \item \(0\)
    \item \(1\)
    \item \(-3\)
  \end{itemize}
  \pause
  \begin{block}{答案}
    \(-2\)；\quad \(-3\)；\quad \(1\)。\quad（都等于 \(-a-2\)）
  \end{block}
\end{frame}

\begin{frame}{四条都齐了}
  \begin{center}
  \begin{tabular}{@{}lcccc@{}}
    \toprule
     & 封闭 & 结合 & 单位元 & 逆元 \\
    \midrule
    模 \(12\) 加法 & \(\checkmark\) & \(\checkmark\) & \(0\) & \(12-a\) \\
    有理数加法 & \(\checkmark\) & \(\checkmark\) & \(0\) & \(-a\) \\
    有理数乘法 & \(\checkmark\) & \(\checkmark\) & \(1\) & 缺 \(0\) 的 \\
    \(a\oplus b=a+b+1\) & \(\checkmark\) & \(\checkmark\) & \(-1\) & \(-a-2\) \\
    \bottomrule
  \end{tabular}
  \end{center}
  \pause
  \begin{center}
    {\large 换了单位元，换了逆元，四条规则照样能凑齐。}
  \end{center}
\end{frame}

% ============================================================================
\section{结合律不是废话}

\begin{frame}{结合律不是废话}
  再换一个运算：
  \[
    a\oplus b=\frac{a+b}{2} .
  \]
  \pause
  算 \((1\oplus3)\oplus5\) 和 \(1\oplus(3\oplus5)\)：
  \pause
  \begin{align*}
    (1\oplus3)\oplus5 &= 2\oplus5 = \frac{2+5}{2} = 3.5,\\
    1\oplus(3\oplus5) &= 1\oplus4 = \frac{1+4}{2} = 2.5 .
  \end{align*}
  \pause
  \begin{att}
    结果不一样。所以这个运算\textbf{不满足结合律}，凑不成群。
  \end{att}
\end{frame}

% ============================================================================
\section{交换律不是必要的}

\begin{frame}{交换律不是必要的}
  看两个动作：
  \begin{block}{}
    动作甲：把数\textbf{乘 \(2\)}。\qquad 动作乙：把数\textbf{加 \(1\)}。
  \end{block}
  \pause
  取 \(x=1\)，先甲后乙：
  \[
    1 \xrightarrow{\ \text{甲}\ } 2 \xrightarrow{\ \text{乙}\ } 3 .
  \]
  \pause
  先乙后甲：
  \[
    1 \xrightarrow{\ \text{乙}\ } 2 \xrightarrow{\ \text{甲}\ } 4 .
  \]
  \pause
  \begin{center}
    {\large \(3\neq4\)：先做哪个，结果不一样。}
  \end{center}
\end{frame}

\begin{frame}{但结合律还在}
  三个动作连着做：甲、乙、甲。
  \begin{itemize}
    \item 先合前两个：把（甲，乙）看成一个动作，再接着做甲；
    \item 先合后两个：先做甲，再把（乙，甲）看成一个动作。
  \end{itemize}
  \pause
  \begin{center}
    {\large 两条路都是「甲、乙、甲」连着做，结果当然一样。}
  \end{center}
  \pause
  \begin{att}
    两件事都记下来：\\
    先做甲、还是先做乙，结果不一样（\(1\to2\to3\) 与 \(1\to2\to4\)）；\\
    而三个动作连着做，先做前两个、还是先做后两个，做的事完全一样。
  \end{att}
  \pause
  \begin{center}
    {\large 这两件事，我们下节课看三角形的时候会再看一遍——\\
    那时候我们能把这组动作数清楚，也能真的验一验。}
  \end{center}
\end{frame}

\begin{frame}{这节课要带走的三句话}
  \begin{itemize}
    \item 群 \(=\) 一个集合 \(+\) 一个运算 \(+\) 四条规则：封闭、结合、单位元、逆元。
    \item 单位元不一定是 \(0\)——在 \(a\oplus b=a+b+1\) 里它是 \(-1\)。
    \item 交换律\textbf{不在}群的规则里。
  \end{itemize}
  \pause
  \begin{center}
    {\large\textcolor{textgreen}{\bfseries 下节课：去看三角形。它的动作也不交换。}}
  \end{center}
\end{frame}

% ---------------- 结尾页 ----------------
\begin{frame}
  \begin{center}
    \vspace{1.6cm}
    {\Large\textcolor{textgreen}{\bfseries 下课}}

    \vspace{1.0cm}
    {\large\textcolor{accentgreen}{下节课：三角形的对称与置换记号}}
  \end{center}
\end{frame}

\end{document}